% lunbian-de-hunling.tex - \usepackage{luatex-cn} 横排现代中文示例
%
% 鲁迅《論辯的魂靈》，一九二五年三月九日刊《語絲》週刊，後收入《華蓋集》。
% 文本取自维基文库（作者 1936 年逝世，属公有领域）。
%
% 本文嵌套引号（「…『…』…」）、叹问连用、省略号、长段引语密集，
% 是横排 clreq 管道的天然试金石：台湾式居中标点、连续标点缩减（2→1.5 字宽）、
% 行首行尾标点挤压、行首行尾禁则、段末孤字避免全部自然生效。
%
% 编译：OSFONTDIR 指向 test/fonts（或安装 TW-Kai）后 lualatex 本文件
\documentclass[12pt]{article}
\usepackage[letterpaper,
            top=3cm, bottom=2.8cm, left=2.75cm, right=2.75cm]{geometry}
\usepackage{fontspec}
% 正文：思源宋体 TC（台湾式居中标点、字面疏朗）；标题保留 TW-Kai 楷体
\setmainfont{Source Han Serif TC}[Script=CJK, Language={Chinese Traditional}]
\newfontfamily\kaiti{TW-Kai}
% 简体部分：思源宋体 SC 正文 + 文鼎PL简中楷（1999 年文鼎科技开源的印刷体楷书）
\newfontfamily\scserif{Source Han Serif SC}[Script=CJK, Language={Chinese Simplified}]
\newfontfamily\kaitisc{AR PL KaitiM GB}
% 说明框内的西文用 EB Garamond（TeX Live 自带），与楷体搭配
\newfontfamily\garamond{EB Garamond}[Scale=1.06]
\newcommand\en[1]{{\garamond #1}}
\usepackage[style=taiwan]{luatex-cn}
\usepackage{xcolor}
\linespread{1.6}
\pagestyle{empty}
\begin{document}

{\centering\kaiti\字号{二号} 論辯的魂靈\par}
\vspace{0.3em}
{\centering\kaiti\字号{小四} 魯迅\par}
\vspace{1.2em}

% 排版说明：楷体 + 米黄底纹的提示框样式，与正文宋体区分（本节为繁体，
% 说明文字亦用繁体）
\noindent\setlength{\fboxsep}{0.8em}%
\colorbox{yellow!14!orange!6}{\parbox{\dimexpr\linewidth-2\fboxsep\relax}{%
  \kaiti\small\setlength{\parindent}{0pt}\color{black!75}%
  【排版說明】正文採用思源宋體 \en{TC}，標題與本框採用 \en{TW-Kai} 楷體，西文採用 \en{EB Garamond}；採用 \en{luatex-cn} 橫向排版，並使用 \en{style=taiwan}（臺灣式置中標點）；其餘採用預設設定──基本級禁則、連續標點縮減一點五字寬、行首與行尾標點擠壓、避免段末孤字；文末落款採右對齊（\en{last-line} 選項）。}}
\vspace{0.8em}

二十年前到黑市，買得一張符，名叫「鬼畫符」。雖然不過一團糟，但帖在壁上看起來，卻隨時顯出各樣的文字，是處世的寶訓，立身的金箴。今年又到黑市去，又買得一張符，也是「鬼畫符」。但帖了起來看，也還是那一張，並不見什麼增補和修改。今夜看出來的大題目是「論辯的魂靈」；細注道：「祖傳老年中年青年『邏輯』扶乩滅洋必勝妙法太上老君急急如律令敕」。今謹摘錄數條，以公同好——

「洋奴會說洋話。你主張讀洋書，就是洋奴，人格破產了！受人格破產的洋奴崇拜的洋書，其價值從可知矣！但我讀洋文是學校的課程，是政府的功令，反對者，即反對政府也。無父無君之無政府党，人人得而誅之。」

「你說中國不好。你是外國人麼？為什麼不到外國去？可惜外國人看你不起……。」

「你說甲生瘡。甲是中國人，你就是說中國人生瘡了。既然中國人生瘡，你是中國人，就是你也生瘡了。你既然也生瘡，你就和甲一樣。而你只說甲生瘡，則竟無自知之明，你的話還有什麼價值？倘你沒有生瘡，是說誑也。賣國賊是說誑的，所以你是賣國賊。我罵賣國賊，所以我是愛國者。愛國者的話是最有價值的，所以我的話是不錯的，我的話既然不錯，你就是賣國賊無疑了！」

「自由結婚未免太過激了。其實，我也並非老頑固，中國提倡女學的還是我第一個。但他們卻太趨極端了，太趨極端，即有亡國之禍，所以氣得我偏要說『男女授受不親』。況且，凡事不可過激；過激派都主張共妻主義的。乙贊成自由結婚，不就是主張共妻主義麼？他既然主張共妻主義，就應該先將他的妻拿出來給我們『共』。」

「丙講革命是為的要圖利：不為圖利，為什麼要講革命？我親眼看見他三千七百九十一箱半的現金抬進門。你說不然，反對我麼？那麼，你就是他的同黨。嗚呼，黨同伐異之風，於今為烈，提倡歐化者不得辭其咎矣！」

「丁犧牲了性命，乃是鬧得一塌糊塗，活不下去了的緣故。現在妄稱志士，諸君切勿為其所愚。況且，中國不是更壞了麼？」

「戊能算什麼英雄呢？聽說，一聲爆竹，他也會吃驚。還怕爆竹，能聽槍炮聲麼？怕聽槍炮聲，打起仗來不要逃跑麼？打起仗來就逃跑的反稱英雄，所以中國糟透了。」

「你自以為是『人』，我卻以為非也。我是畜類，現在我就叫你爹爹。你既然是畜類的爹爹，當然也就是畜類了。」

「勿用驚歎符號，這是足以亡國的。但我所用的幾個在例外。

中庸太太提起筆來，取精神文明精髓，作明哲保身大吉大利格言二句云：

中學為體西學用，

不薄今人愛古人。」

\vspace{0.6em}
\horiSetup{last-line=right}
{\small 一九二五年三月九日刊《語絲》週刊第十七期，後收入《華蓋集》。

文本取自維基文庫，屬公有領域。}

\horiSetup{last-line=left}

% ============================================================================
% 简体·中国式排版（同一篇文章的对照排法）
% ============================================================================
\clearpage
\scserif
\horiSetup{style=mainland, quote-style=curly}

{\centering\kaitisc\字号{二号} lunbian-de-hunling\par}
\vspace{0.3em}
{\centering\kaitisc\字号{小四} 鲁迅\par}
\vspace{1.2em}

\noindent\setlength{\fboxsep}{0.8em}%
\colorbox{blue!7!cyan!5}{\parbox{\dimexpr\linewidth-2\fboxsep\relax}{%
  \kaitisc\small\setlength{\parindent}{0pt}\color{black!75}%
  【排版说明】同一篇文章的中国式排法：正文用思源宋体 \en{SC}，标题与本框用文鼎简中楷；\en{style=mainland}（标点偏靠始端、空白在末端），并启用 \en{quote-style=curly}——原文的直角引号自动转为弯引号，嵌套体例先双后单。}}
\vspace{0.8em}

二十年前到黑市，买得一张符，名叫「鬼画符」。虽然不过一团糟，但帖在壁上看起来，却随时显出各样的文字，是处世的宝训，立身的金箴。今年又到黑市去，又买得一张符，也是「鬼画符」。但帖了起来看，也还是那一张，并不见什么增补和修改。今夜看出来的大题目是「lunbian-de-hunling」；细注道：「祖传老年中年青年『逻辑』扶乩灭洋必胜妙法太上老君急急如律令敕」。今谨摘录数条，以公同好——

「洋奴会说洋话。你主张读洋书，就是洋奴，人格破产了！受人格破产的洋奴崇拜的洋书，其价值从可知矣！但我读洋文是学校的课程，是政府的功令，反对者，即反对政府也。无父无君之无政府党，人人得而诛之。」

「你说中国不好。你是外国人么？为什么不到外国去？可惜外国人看你不起……。」

「你说甲生疮。甲是中国人，你就是说中国人生疮了。既然中国人生疮，你是中国人，就是你也生疮了。你既然也生疮，你就和甲一样。而你只说甲生疮，则竟无自知之明，你的话还有什么价值？倘你没有生疮，是说诳也。卖国贼是说诳的，所以你是卖国贼。我骂卖国贼，所以我是爱国者。爱国者的话是最有价值的，所以我的话是不错的，我的话既然不错，你就是卖国贼无疑了！」

「自由结婚未免太过激了。其实，我也并非老顽固，中国提倡女学的还是我第一个。但他们却太趋极端了，太趋极端，即有亡国之祸，所以气得我偏要说『男女授受不亲』。况且，凡事不可过激；过激派都主张共妻主义的。乙赞成自由结婚，不就是主张共妻主义么？他既然主张共妻主义，就应该先将他的妻拿出来给我们『共』。」

「丙讲革命是为的要图利：不为图利，为什么要讲革命？我亲眼看见他三千七百九十一箱半的现金抬进门。你说不然，反对我么？那么，你就是他的同党。呜呼，党同伐异之风，于今为烈，提倡欧化者不得辞其咎矣！」

「丁牺牲了性命，乃是闹得一塌糊涂，活不下去了的缘故。现在妄称志士，诸君切勿为其所愚。况且，中国不是更坏了么？」

「戊能算什么英雄呢？听说，一声爆竹，他也会吃惊。还怕爆竹，能听枪炮声么？怕听枪炮声，打起仗来不要逃跑么？打起仗来就逃跑的反称英雄，所以中国糟透了。」

「你自以为是『人』，我却以为非也。我是畜类，现在我就叫你爹爹。你既然是畜类的爹爹，当然也就是畜类了。」

「勿用惊叹符号，这是足以亡国的。但我所用的几个在例外。

中庸太太提起笔来，取精神文明精髓，作明哲保身大吉大利格言二句云：

中学为体西学用，

不薄今人爱古人。」

\vspace{0.6em}
\horiSetup{last-line=right}
{\small 一九二五年三月九日刊《语丝》周刊第十七期，后收入《华盖集》。

简体文本据维基文库繁体原文转换，属公有领域。}

\horiSetup{last-line=left}

\end{document}
